%\documentclass[11pt,a4paper,twoside]{frontmatter/tesis}
% SI NO PENSAS IMPRIMIRLO EN FORMATO LIBRO PODES USAR
\documentclass[11pt,a4paper]{tesis}

% subfiles structure
%\usepackage{subfiles}
\usepackage[subpreambles=true]{standalone}
\usepackage{import}

%other packages
\usepackage{graphicx}
\usepackage[utf8]{inputenc}
\usepackage[spanish]{babel}
\usepackage[left=3cm,right=3cm,bottom=3.5cm,top=3.5cm]{geometry}

%\bibliographystyle{apalike}
\usepackage[authoryear,round]{natbib}

\bibliographystyle{myplainnat}


\renewcommand{\betweenauthors}{&}
%\usepackage[style=authoryear]{biblatex}

%general adjustments
%\setcounter{tocdepth}{3}
%\graphicspath{ {images/}{../../images}] }


%librerias chapter 1
\usepackage{amssymb}
\usepackage{amsmath}
%\usepackage{natbib}
\usepackage{subfigure}
\usepackage{gensymb}
\usepackage{authblk}
\usepackage{url}


\usepackage{hyperref}
\usepackage{cleveref}
\usepackage{nameref}

\usepackage[subfigure]{tocloft}


%para hacer comentarios
\usepackage{xargs}                      % Use more than one optional parameter in a new commands
\usepackage[pdftex,dvipsnames]{xcolor}  % Coloured text etc.
\usepackage[colorinlistoftodos,prependcaption,textsize=tiny]{todonotes}
\newcommandx{\unsure}[2][1=]{\todo[linecolor=red,backgroundcolor=red!25,bordercolor=red,#1]{#2}}
\newcommandx{\change}[2][1=]{\todo[linecolor=blue,backgroundcolor=blue!25,bordercolor=blue,#1]{#2}}
\newcommandx{\pendiente}[2][1=]{\todo[linecolor=OliveGreen,backgroundcolor=OliveGreen!25,bordercolor=OliveGreen,#1]{#2}}
\newcommandx{\info}[2][1=]{\todo[linecolor=Plum,backgroundcolor=Plum!25,bordercolor=Plum,#1]{#2}}
\newcommandx{\thiswillnotshow}[2][1=]{\todo[disable,#1]{#2}}


%%%%%%%%%%%%% Titulos

%%%%%%%%%%%% paragraph con titulo en negrita
\newcommand{\partitle}[1]{%
	\addvspace{0.2cm}% perhaps, as Barbara Beeton suggests
	\textbf{#1}
	\par
	\hspace{0.5cm}
}


%%%%%%%%% \subsubsection en negrita e italica
\usepackage{titlesec}

\titleformat{\subsubsection}
{\normalfont\fontfamily{phv}\fontsize{10}{17}\bfseries\itshape}{\thesubsection}{1em}{}



\begin{document}

%%%% CARATULA

\def\autor{Diego Kozlowski}
\def\tituloTesis{Técnicas de Ciencia de Datos aplicadas al comercio internacional \change{provisorio}}
\def\runtitulo{Técnicas de Ciencia de Datos aplicadas al comercio internacional \change{provisorio}}
\def\runtitle{Data Science Techniques applied to World Trade \change{provisorio}}
\def\director{Dra. Viktoriya Semeshenko}
%\def\codirector{}
\def\lugar{Buenos Aires, 2019}
\newcommand{\HRule}{\rule{\linewidth}{0.2mm}}
%
\thispagestyle{empty}

\begin{center}\leavevmode

\vspace{-2cm}

\begin{tabular}{l}
\includegraphics[width=2.6cm]{frontmatter/logofcen.pdf}
\end{tabular}


{\large \sc Universidad de Buenos Aires

Facultad de Ciencias Exactas y Naturales

Maestría en Explotación de Datos y Descubrimiento del Conocimiento}

\vspace{2.0cm}

%\vspace{3.0cm}
%{
%\Large \color{red}
%\begin{tabular}{|p{2cm}cp{2cm}|}
%\hline
%& Pre-Final Version: \today &\\
%\hline
%\end{tabular}
%}
%\vspace{2.5cm}

\begin{huge}
\textbf{\tituloTesis}
\end{huge}

\vspace{2cm}

{\large Tesis presentada para optar al Título de Magister de la Universidad de Buenos Aires en Explotación de Datos y Descubrimiento del Conocimiento}

\vspace{2cm}

{\Large \autor}

\end{center}

\vfill

{\large

{Directora: \director}

\vspace{1cm}

%{Codirector: \codirector}


\lugar
\vspace{1cm}

\fechapre

\vspace{1cm}
\fechadef

}

\newpage\thispagestyle{empty}


%%%% ABSTRACTS, AGRADECIMIENTOS Y DEDICATORIA
\frontmatter
\pagestyle{empty}
%\begin{center}
%\large \bf \runtitulo
%\end{center}
%\vspace{1cm}
\chapter*{\runtitulo}

\noindent La princesa Leia, líder del movimiento rebelde que desea reinstaurar la República en la galaxia en los tiempos ominosos del Imperio, es capturada por las malévolas Fuerzas Imperiales, capitaneadas por el implacable Darth Vader. El intrépido Luke Skywalker, ayudado por Han Solo, capitán de la nave espacial ``El Halcón Milenario'', y los androides, R2D2 y C3PO, serán los encargados de luchar contra el enemigo y rescatar a la princesa para volver a instaurar la justicia en el seno de la Galaxia (aprox. 200 palabras).

\bigskip

\noindent\textbf{Palabras claves:} Guerra, Rebelión, Wookie, Jedi, Fuerza, Imperio (no menos de 5).

\cleardoublepage
%\begin{center}
%\large \bf \runtitle
%\end{center}
%\vspace{1cm}
\chapter*{\runtitle}

\noindent In a galaxy far, far away, a psychopathic emperor and his most trusted servant -- a former Jedi Knight known as Darth Vader -- are ruling a universe with fear. They have built a horrifying weapon known as the Death Star, a giant battle station capable of annihilating a world in less than a second. When the Death Star's master plans are captured by the fledgling Rebel Alliance, Vader starts a pursuit of the ship carrying them. A young dissident Senator, Leia Organa, is aboard the ship \& puts the plans into a maintenance robot named R2-D2. Although she is captured, the Death Star plans cannot be found, as R2 \& his companion, a tall robot named C-3PO, have escaped to the desert world of Tatooine below. Through a series of mishaps, the robots end up in the hands of a farm boy named Luke Skywalker, who lives with his Uncle Owen \& Aunt Beru. Owen \& Beru are viciously murdered by the Empire's stormtroopers who are trying to recover the plans, and Luke \& the robots meet with former Jedi Knight Obi-Wan Kenobi to try to return the plans to Leia Organa's home, Alderaan. After contracting a pilot named Han Solo \& his Wookiee companion Chewbacca, they escape an Imperial blockade. But when they reach Alderaan's coordinates, they find it destroyed - by the Death Star. They soon find themselves caught in a tractor beam \& pulled into the Death Star. Although they rescue Leia Organa from the Death Star after a series of narrow escapes, Kenobi becomes one with the Force after being killed by his former pupil - Darth Vader. They reach the Alliance's base on Yavin's fourth moon, but the Imperials are in hot pursuit with the Death Star, and plan to annihilate the Rebel base. The Rebels must quickly find a way to eliminate the Death Star before it destroys them as it did Alderaan (aprox. 200 palabras).

\bigskip

\noindent\textbf{Keywords:} War, Rebellion, Wookie, Jedi, The Force, Empire (no menos de 5). % OPCIONAL: comentar si no se quiere

\cleardoublepage
\chapter*{Agradecimientos}

\noindent Lorem ipsum dolor sit amet, consectetur adipiscing elit. Fusce sapien ipsum, aliquet eget convallis at, adipiscing non odio. Donec porttitor tincidunt cursus. In tellus dui, varius sed scelerisque faucibus, sagittis non magna. Vestibulum ante ipsum primis in faucibus orci luctus et ultrices posuere cubilia Curae; Mauris et luctus justo. Class aptent taciti sociosqu ad litora torquent per conubia nostra, per inceptos himenaeos. Mauris sit amet purus massa, sed sodales justo. Mauris id mi sed orci porttitor dictum. Donec vitae mi non leo consectetur tempus vel et sapien. Curabitur enim quam, sollicitudin id iaculis id, congue euismod diam. Sed in eros nec urna lacinia porttitor ut vitae nulla. Ut mattis, erat et laoreet feugiat, lacus urna hendrerit nisi, at tincidunt dui justo at felis. Class aptent taciti sociosqu ad litora torquent per conubia nostra, per inceptos himenaeos. Ut iaculis euismod magna et consequat. Mauris eu augue in ipsum elementum dictum. Sed accumsan, velit vel vehicula dignissim, nibh tellus consequat metus, vel fringilla neque dolor in dolor. Aliquam ac justo ut lectus iaculis pharetra vitae sed turpis. Aliquam pulvinar lorem vel ipsum auctor et hendrerit nisl molestie. Donec id felis nec ante placerat vehicula. Sed lacus risus, aliquet vel facilisis eu, placerat vitae augue.
 % OPCIONAL: comentar si no se quiere

%\cleardoublepage
%\hfill \textit{A mi persona favorita.}
  % OPCIONAL: comentar si no se quiere

\cleardoublepage
\tableofcontents

\cleardoublepage
\listoffigures

\mainmatter
\pagestyle{headings}

%%%% ACA VA EL CONTENIDO DE LA TESIS

\chapter{introducción}\label{sec:introduccion}

\import{chapters/}{chapter_1/chapter1}

\chapter{Comercio a nivel países}\label{sec:agregado}

\import{chapters/}{chapter_2/chapter2}

%\section{Introducción}
%\info{inclyue estado del arte}
%
%\section{Análisis exploratorio de datos}
%
%\section{Metodología}
%
%\section{Resultados}
%
%\subsection{El comercio internacional en el 2016 a la luz del análisis de redes}
%
%

%---------------------------

\chapter[Comercio a nivel producto]{Comercio a nivel producto} \label{sec:desagregado}




\import{chapters/}{chapter_3/chapter3}



\chapter{Conclusiones} \label{sec:conclusiones}


\import{chapters/}{chapter_4/chapter4}

%\appendix 
\chapter*{Anexo}\label{append}
\addcontentsline{toc}{section}{Anexo}
\renewcommand{\thesubsection}{\Alph{subsection}}

\import{chapters/}{Anexo/anexo}



%%%% BIBLIOGRAFIA
\backmatter
%\bibliographystyle{plainnat}
\bibliography{backmatter/bibliografia}

%\bibliography{bibliografia}

\end{document}
